% Chapter One: the problem, the argument and the plan of the dissertation.
\chapter{Introduction}
\label{ch:intro}

In the spring of 1412 a Lüneburg carter named Hinrik Vos paid four pennies at the toll bridge of Mölln for a wagon of salt bound for Lübeck. The entry is unremarkable; the same register records more than three hundred such payments that year. Yet entries of this kind, read in aggregate, describe something that the familiar narrative of the Hanseatic League has largely ignored: a web of overland routes, carriers and toll stations without which the maritime trade of the Baltic could not have functioned.

This dissertation reconstructs that web. It asks how goods moved between the inland centres of production and the coastal ports of the southern Baltic between 1350 and 1450, who controlled the points at which routes converged, and how the network responded to the demographic shocks of the late fourteenth century.

\section{The Maritime Bias of Hanseatic History}

The Hanse has been written about chiefly as a league of seaports. Its great nineteenth-century historians took their cue from the surviving records of the \emph{Hansetage}, the diets at which delegates of the member towns negotiated privileges abroad, and these records are overwhelmingly concerned with ships, cargoes and foreign counting houses \parencite{dollinger1964,hammel1988}. The resulting picture is one of a fleet with towns attached.\footnote{The phrase is \textcite{brandt1997}'s, who uses it to describe, and to criticise, the older consensus.}

More recent work has complicated this picture without overturning it. Studies of individual merchant firms show that the same families who freighted ships at Lübeck also financed wagon trains to Nuremberg and Kraków \parencite{selzer2003}. Studies of the salt trade show that Lüneburg salt reached Scandinavia only after a journey of some eighty kilometres overland \parencite{witthoft1976}. But these findings remain scattered, and no study has yet treated the land routes as a system in their own right.

\section{Argument}

The central claim of this dissertation is that the overland routes of the southern Baltic formed a coherent network whose structure, rather than the volume of goods it carried, determined which towns prospered. Three subsidiary claims support it.

First, the network was \emph{sparse but well connected}: most towns traded with only a few partners, yet any two towns could be linked through a short chain of intermediaries. Second, a small number of inland toll towns occupied brokerage positions, sitting on the only paths between otherwise separate regional systems. Third, the plague years did not destroy the network but reorganised it, concentrating traffic on fewer, better-maintained roads and thereby strengthening the brokers.

\section{Plan of the Dissertation}

Chapter~\ref{ch:sources} introduces the sources, chiefly eleven surviving toll registers, and the procedures used to make their entries comparable. It also sets out the two quantitative tools used throughout: a gravity model of bilateral trade and a measure of brokerage drawn from social network analysis. Chapter~\ref{ch:salt} applies these tools to the salt roads that connected Lüneburg to the Baltic coast, tracing how traffic and brokerage shifted across the century. Appendix~\ref{app:registers} lists the registers and summarises their coverage.
